\chapter{Introduction}\label{chap: Introduction}

\fxchange{Vllt etwas allgemeiner beginnen: Vorschlag von ChatGPT}
The principles of fluid mechanics are essential in numerous disciplines, including mechanical, civil, aerospace, and chemical engineering. By applying the laws of conservation of mass, momentum, and energy, engineers can predict and analyze complex flow phenomena. 2 Zitate
 The study of fluid flow enables the design and optimization of systems such as pipelines, pumps, turbines, aircraft, and hydraulic structures. Mathematical modeling and computational methods have significantly enhanced the ability to simulate and evaluate fluid behavior under various operating conditions. 
 
 Turbulence is a complex flow regime characterized by chaotic and irregular fluctuations in velocity, pressure, and other flow properties.  
%Turbulent flows are present in almost all engineering applications.
In particular,  turbulent flows over a wall, so known as wall-bounded turbulent flows, are  frequently encountered  in various kinds of technical applications ZITAT. \fxchange{gerne 2-3 Anwendungen nennen}
Wall-bounded turbulence is of great interest in research and technology since around $50 \%$ of the energy consumption in pipe flows  or the movement of cars or planes through the air can be attributed to the dissipation of energy in wall-bounded turbulence \citep{jimenezCascadesWallBoundedTurbulence2012}.
In order to reduce this figure, an better understanding of fluid dynamics in general, and  of turbulence behaviour at the wall and in the near-wall region in particular, is essential.

An analytical solution to the governing equations of fluid dynamics, as  discussed in \autoref{chap: Theory}, is impossible in most cases.
Apart from experimental setups, the use of \acrfull{cfd} is a vital tool in fluid dynamics research to increase the understanding.
\acrfull{dns} and \acrfull{les} are particularly popular among  researchers, since they can provide valuable insights into  turbulent flow behaviour \citep{moinDIRECTNUMERICALSIMULATION1998}.
If turbulent structures do not need to be resolved or the design and optimisation of engineering tasks are the goal, \acrfull{rans} simulations are a common choice involving significantly less computational effort.


In many real-world applications,  such as shipping and aviation, wall roughness has a major impact on drag.
In shipping, the biofouling-induced roughness can increase the overall drag by approximately $10 \%$ \citep{flackReviewHydraulicRoughness2010}. 
In aviation, even the comparatively low roughness on aircraft surfaces accounts for approximately $5.5 \%$ of drag \citep{chungPredictingDragRough2021}. 

Many studies have been conducted in order to improve the understanding of turbulence behaviour in rough walled conditions.
One notable early experimental study was carried out by~\citeauthor{StroemungsgesetzeNikuradse}, who investigated the behaviour of turbulent flows in rough pipes.
The roughness was created by attaching sand grains  to the pipe wall \citep{StroemungsgesetzeNikuradse}. 
Based on  the data of~\citeauthor{StroemungsgesetzeNikuradse}, the concept of the equivalent sand grain roughness was developed by~\cite{schlichtingExperimentelleUntersuchungenRauhigkeitsproblem1936}.
More recent experimental efforts were carried out by~\cite{schultzRoughwallTurbulentBoundary2007a} and~\cite{frohnapfelFlowResistanceHeterogeneous2024} who used sandpaper as the roughness geometry.
Experimental studies of generated irregular rough surfaces have also been conducted, e.g., in \citep{flackSkinFrictionMeasurements2020}, in which the surfaces were generated using an algorithm and produced by a 3D-printer.


Substantial efforts were also made to investigate the behaviour of turbulence over rough surfaces using simulations.
Since the behaviour of the turbulence is of interest, almost all studies were conducted using the \acrshort{dns} method.
Turbulent flows over rough surfaces are mostly studied in the canonical channel case, e.g.\ by \citep{thakkarInvestigationTurbulentFlow2017,flackSkinFrictionMeasurements2020,kimTurbulenceStatisticsFully1987,jellyReynoldsDispersiveShear2018}, while the turbulent rough pipe has received less attention \citep{chanSystematicInvestigationRoughness2015,demaioDirectNumericalSimulation2023}.
Some studies have used isolated roughness elements to determine the impact of individual roughness features \citep{chatzikyriakouDNSTurbulentFlow2015}, while others have attempted to investigate the real-world turbulence behaviour by using scanned surfaces of real objects \citep{thakkarSurfaceCorrelationsHydrodynamic2017,flackSkinFrictionMeasurements2020,frohnapfelFlowResistanceHeterogeneous2024}.
Furthermore, both regular and irregular roughness have been investigated.
Examples of regular roughness include the work of~\citeauthor{leonardiStructureTurbulentChannel2004} who placed square bars onto a flat surface, and~\citeauthor{chanSystematicInvestigationRoughness2015} who investigated turbulent pipe flow with a sinusoidal surface \citep{chanSystematicInvestigationRoughness2015,chanSecondaryMotionTurbulent2018,leonardiStructureTurbulentChannel2004}.
Irregular roughness is often studied using generated irregular surfaces, e.g.\ \cite{thakkarInvestigationTurbulentFlow2017,busseReynoldsnumberDependenceNearwall2017} and~\cite{demaioDirectNumericalSimulation2023}.

Despite these numerous efforts, no general law yet exists to predict flow characteristics, e.g.\ drag or roughness functions, directly from the surface topography \citep{chungPredictingDragRough2021}.
Furthermore, the past \acrshort{dns} studies mostly used the \acrfull{ibm} to achieve rough surfaces in their simulation setups, whereas a body-fitted mesh is used in the present thesis.\fxchange{Dieser Halbsatz gehört erst in den nächsten Absatz}

The aim of this thesis is to provide additional data to enhance the understanding  of turbulent pipe flow behaviour over rough surfaces.
%Additionally, a body-fitted mesh topology is chosen deliberately, to provide date for an 
Additionally, a body-fitted mesh topology has been deliberately chosen to provide data for \acrshort{dns} studies in rough pipes, which are performed using the Neko solver.
 ZITAT
The influence of changing Reynolds numbers and surface parameters is investigated using the results of the simulations.\fxchange{hier dürfen gerne noch 5 weitere Sätze zum methodischen Vorgehen folgen insb darfst Du dich von der bisher zitierten Liteartur abgrenzen!}
These results are documented in this thesis.\fxchange{das darf man gerne mit den 3 wichtigsten Erkenntnissen auswalzen}

\fxchange{hier gehört ein neuer Absatz hin}
\section{Structure of the thesis}
\autoref{chap: Introduction} provides motivation for the relevance of the topic.
Furthermore, the  literature is reviewed briefly and the topic of the thesis is explained.
\autoref{chap: Theory} explains the theoretical background to wall-bounded turbulent flow and provides an overview  of the influence of wall roughness on turbulent flows.
\autoref{chap: Methods} provides an explanation of the methods used.
\acrshort{dns} and the \acrshort{sem} are introduced.
Additionally, the algorithm used to generate the rough surfaces is briefly discussed.
\autoref{chap: application} presents the simulation setup used to perform the \acrshort{dns}, followed by a description of  how roughness is applied onto the pipe mesh.
The post-processing of rough surfaces is also explained, as are  the  simulations performed.
In \autoref{chap: Validation}, the simulation setup is tested on smooth pipe geometries and the data is compared with that in the  literature to validate the methodology.
Furthermore, the convergence studies  conducted  to find reasonable simulation parameters are presented.
\autoref{chap: Results} discusses the results of the simulations performed.
Plots of the flow statistics, as well as instantaneous fields and energy spectra are provided.
The different setups are compared and interpretations of the results are given.
Finally, \autoref{chap: Conclusion} concludes the thesis and provides an outlook for potential future work.


